% !PRINTED 283
\weberpage{283}

Führen wir ein zweites System von Variablen $\tau$ ein und setzen
\[
z = \tau_{n-1} f_0 + \tau_{n-2} f_1 + \cdots + \tau_0 f_{n-1},
\]
so ergiebt sich aus (4) mit Benutzung der Formeln [\S.\,68, (6)]
\[
S(f_s) = (n-s)\,a_s
\]
\weq{6}{\begin{aligned}
S[(x-\alpha)\,y\,z] &= n\,a_0 \sum_{0,\,n-1}^{s} (E_{-1,s} - \alpha\,E_{0,s})\,\tau_{n-s-1} \\
&\quad + (n-1)\,a_1 \sum_{0,\,n-1}^{s} (E_{0,s} - \alpha\,E_{1,s})\,\tau_{n-s-1} \\
&\quad \cdots\cdots\cdots\cdots\cdots\cdots\cdots \\
&\quad + a_{n-1} \sum_{0,\,n-1}^{s} (E_{n-2,s} - \alpha\,E_{n-1,s})\,\tau_{n-s-1},
\end{aligned}}
und dies geht geradezu in die Function $H$ über, wenn wir $\tau = t$ setzen.

Nehmen wir als Beispiel den Fall $n = 3$, so ergiebt sich
\[
\begin{aligned}
E_{-1,0} &= -a_1 t_2 - a_2 t_1 - a_3 t_0, & E_{0,0} &= a_0 t_2, \\
E_{-1,1} &= -a_2 t_2 - a_3 t_1, & E_{0,1} &= -a_2 t_1 - a_3 t_0, \\
E_{-1,2} &= -a_3 t_2, & E_{0,2} &= -a_3 t_1, \\
E_{1,0} &= a_0 t_1, & E_{2,0} &= a_0 t_0, \\
E_{1,1} &= a_0 t_2 + a_1 t_1, & E_{2,1} &= a_0 t_1 + a_1 t_0, \\
E_{1,2} &= -a_3 t_0, & E_{2,2} &= a_0 t_2 + a_1 t_1 + a_2 t_0,
\end{aligned}
\]
und hieraus kann man die Coëfficienten der Form $H$ nach (6) leicht berechnen:
\[
\begin{aligned}
H_{2,2} &= -a_0 (a_1 + 3 a_0 \alpha) \\
H_{1,1} &= -3 a_0 a_3 - a_1 a_2 - 2(a_1^2 - a_0 a_2)\alpha \\
H_{0,0} &= -a_2 a_3 + (2 a_1 a_3 - a_2^2)\alpha \\
H_{1,0} &= -2 a_1 a_3 + (3 a_0 a_3 - a_1 a_2)\alpha \\
H_{2,1} &= -2 a_0 (a_2 + a_1 \alpha) \\
H_{0,2} &= -a_0 (3 a_3 + a_2 \alpha).
\end{aligned}
\]

\begin{center}
\S.\,92.

\gesperrt{Die Determinante der Hermite'schen Form.}
\end{center}

Für die Frage nach der Anzahl der negativen Quadrate der Form $H$ ist die Kenntniss ihrer Determinante von Wichtigkeit. Diese Determinante lässt sich allgemein auf folgende Art berechnen.
